% 02_inicio_rapido.tex - Manual Capítulo 2: Inicio Rápido
\chapter{Inicio Rápido}
\label{man:inicio}

\section{Ejecutar la Interfaz Gráfica (GUI)}
\label{sec:gui-inicio}

\begin{minted}{bash}
# Abrir la GUI con un escenario específico
uv run python -m tactical_grid.gui scenarios/prueba_20x20.json

# Sin argumentos: abre diálogo para elegir archivo
uv run python -m tactical_grid.gui
\end{minted}

La ventana se abre mostrando:
\begin{itemize}
\item \textbf{Mapa central:} Grid con terrenos coloreados por costo (verde claro = barato, marrón = caro), obstáculos gris oscuro, bases con borde azul/rojo, recurso (diamante dorado), unidades (círculos azul/rojo).
\item \textbf{Panel derecho:} Turno actual, movimiento \#/límite, portador del recurso, unidad seleccionada, meta de búsqueda y métricas.
\item \textbf{Traza:} Orden de expansión de búsquedas, valores de las acciones raíz en decisiones y acciones de la partida.
\item \textbf{Barra inferior:} Botones \textit{Ejecutar}, \textit{Animar}, \textit{Aplicar jugada}, \textit{Reiniciar}, verificación RF-42, selector de bando y \textit{Partida automática}.
\end{itemize}

\section{Ejecutar una Búsqueda (CLI)}
\label{sec:cli-busqueda}

\begin{minted}{bash}
# BFS desde A1 a (3,2) en el ejemplo oficial
uv run python -m tactical_grid scenarios/apendice_a.json \
    --algorithm BFS --unit A1 --goal 3,2 --verify

# A* con heurística obstacle-aware (h2)
uv run python -m tactical_grid scenarios/caso_03_obstaculo_heuristica.json \
    --algorithm A_STAR --heuristic H2 --verify
\end{minted}

La salida es JSON con claves equivalentes al Apéndice B:
\code{algoritmo, exito, camino, costo, longitud, estados\_generados, estados\_expandidos, maximo\_frontera, tiempo\_segundos}.

\section{Jugar una Partida (GUI)}
\label{sec:jugar-gui}

\begin{enumerate}
\item Abre la GUI: \code{uv run python -m tactical\_grid.gui scenarios/prueba\_20x20.json}
\item Modo \textit{Seleccionar / mover unidades} (radio button izquierdo).
\item Clic en una unidad de tu bando (azul = A, rojo = B) $\to$ se selecciona (anillo negro discontinuo).
\item Celdas verdes $\to$ destinos legales para esa unidad. Clic en una $\to$ mueve.
\item Turno pasa automáticamente al otro bando.
\item Gana quien lleva el recurso (diamante dorado) a su base (borde grueso de su color).
\item Empate al llegar al límite de turnos (default 100).
\end{enumerate}

\section{Jugar una Partida desde la Terminal}
\label{sec:partida-cli-rapida}

\begin{minted}{bash}
# IA contra IA hasta terminar la partida
uv run python -m tactical_grid scenarios/caso_07_decision_adversarial.json \
    --mode partida --play auto --algorithm ALPHA_BETA --depth 3

# Persona contra IA; --side B permite jugar segundo
uv run python -m tactical_grid scenarios/caso_07_decision_adversarial.json \
    --mode partida --play humano --side A --depth 3

# Variante dinámica
uv run python -m tactical_grid scenarios/caso_11_dinamico.json \
    --mode partida --play humano --variant dinamica --depth 3
\end{minted}

En el modo \code{humano}, la terminal enumera las acciones legales y se selecciona una
escribiendo su número; \code{q} o \code{salir} abandona la partida.

\section{Atajos de Teclado / Mouse}
\label{sec:atajos}

\begin{table}[H]
\centering
\caption{Atajos en la GUI}
\begin{tabular}{ll}
\toprule
\textbf{Acción} & \textbf{Tecla o clic} \\
\midrule
Seleccionar unidad & Clic izquierdo sobre la unidad \\
Mover unidad & Clic izquierdo en una celda verde (después de seleccionar) \\
Fijar meta de búsqueda & Clic derecho en cualquier celda \\
Aplicar decisión & Botón \textit{Aplicar jugada} \\
Animar búsqueda & Botón \textit{Animar} \\
Partida automática & Botón \textit{Partida automática}; vuelve a pulsarlo para pausar \\
Reiniciar partida & Botón \textit{Reiniciar} \\
\bottomrule
\end{tabular}
\end{table}